\section{Related Work}
\paragraph{Computational psycholinguistics.} There is a rich line of work in computational psycholinguistics on cognitively plausible models of human language processing. Such work focuses chiefly on a) understanding whether high-level parsing methods can be used to predict psycholinguistic data (such as reading-time or eye-tracking data) and neural data (eg. fMRI and ecog data from linguistic experiments) and/or b) developing parsing methods that have specific, hallmark, experimentally-established cognitive properties of human syntactic processing (most importantly {\em incrementality} of parsing, {\em limited memory} constraints, and {\em connectedness} of the syntactic structures maintained by the parser). These goals are summarized in the directional survey paper of Keller \cite{keller}, which summarizes the psycholinguistic desiderata of b) and proposes evaluation standards for a). 

The interested reader should start with Jurafsky’s early work on using probabilistic parsing to predict reading difficulty \cite{jurafsky}, the surprisal-based models of Hale, Levy, Demberg and Keller \cite{Hale2001, Levy2008, DembergKeller2008a}, and the strictly incremental (a “b” goal) predictive models of Demberg and Keller \cite{DembergKellery2008b, DembergKellery2009}.  Much work attempts to achieve the properties of b) while maintaining a)-- neural and psycholinguistic predictiveness;  for instance the PLTAG parsers of Demberg et al \cite{Demberg}, or the parser of Vasisth and Lewis \cite{Vasisth}, which is constructed in a different high-level meta-model of cognitive processing (ACT-R). A related line of work takes modern neural parsing methods, and examines whether the internal states of these models at parse time can be used to model psycholinguistic and neurological data observed for the same sentences-- see \cite{findingsyntax} which uses an neural action-based parser, and \cite{Schimpf} which compares a wide range of ANN methods.

The present work differs from this line of work in key ways (though it is critically related, see end of this section). All the previous work studies parsing in a high-level programming language, whereas we are interested in whether we can implement a simple parser out of individual neurons and synapses, starting with a realistic mathematical model of neuronal processing. Such a simulation would take place on the level of millions of neurons and synapses-- it is non-trivial how to implement even a single step of syntactic processing, whereas in previous work, such structure-building “actions” are built into the framework’s primitives. At present, our work can be seen as largely orthogonal to the related work described above, as we attempt to bridge granular neuronal mechanics with the study of complex cognitive processes such as syntax. See Section [] for discussions of potential future work that connects the two areas.


\paragraph{Neuroscience.} In building a neural parser, we begin with the most established tenets of neuron biology. Individual neurons fire when they receive sufficient excitatory input from other neurons, firing is an atomic operation, synapses between neurons become stronger with repeated co-firing (plasticity), and some synapses can be inhibitory (firing one neuron decreases the total synaptic input to another neuron). These essentials of neuroscience, including plasticity, are covered in a number of textbooks (see [Book ref, Chapters x y z] for an example). 

How are higher-level cognitive processes computed by collections of neurons? Highly interconnected sets of neurons, or assemblies, are an increasingly popular hypothesis for the main unit of higher-level cognition in modern neuroscience. First hypothesized decades ago by Hebb, assemblies have been identified experimentally [cite] and their dynamics have been studied in animal brains [Yuste papers], displacing the previously dominant single-neuron theory of information encoding in the brain (see 
“Barlow versus Hebb: When is it time to abandon the notion of feature detectors and adopt the cell assembly as the unit of cognition?” \cite{eichenbaum}). 

In their recent PNAS paper, Papadimitriou et al. propose a concrete mathematical formalization of assemblies, and demonstrate that this simplified model is capable of simulating basic assembly dynamics that are known experimentally. This model, summarize in section 3, gives us a framework for modeling neurons and assembly dynamics.

\paragraph{Language in the brain.} 


